% Options for packages loaded elsewhere
\PassOptionsToPackage{unicode}{hyperref}
\PassOptionsToPackage{hyphens}{url}
\PassOptionsToPackage{dvipsnames,svgnames,x11names}{xcolor}
%
\documentclass[
  11pt,
  a4paper,
]{ltjsarticle}
\usepackage{amsmath,amssymb}
\usepackage{iftex}
\ifPDFTeX
  \usepackage[T1]{fontenc}
  \usepackage[utf8]{inputenc}
  \usepackage{textcomp} % provide euro and other symbols
\else % if luatex or xetex
  \usepackage{unicode-math} % this also loads fontspec
  \defaultfontfeatures{Scale=MatchLowercase}
  \defaultfontfeatures[\rmfamily]{Ligatures=TeX,Scale=1}
\fi
\usepackage{lmodern}
\ifPDFTeX\else
  % xetex/luatex font selection
\fi
% Use upquote if available, for straight quotes in verbatim environments
\IfFileExists{upquote.sty}{\usepackage{upquote}}{}
\IfFileExists{microtype.sty}{% use microtype if available
  \usepackage[]{microtype}
  \UseMicrotypeSet[protrusion]{basicmath} % disable protrusion for tt fonts
}{}
\makeatletter
\@ifundefined{KOMAClassName}{% if non-KOMA class
  \IfFileExists{parskip.sty}{%
    \usepackage{parskip}
  }{% else
    \setlength{\parindent}{0pt}
    \setlength{\parskip}{6pt plus 2pt minus 1pt}}
}{% if KOMA class
  \KOMAoptions{parskip=half}}
\makeatother
\usepackage{xcolor}
\usepackage[margin=24mm]{geometry}
\setlength{\emergencystretch}{3em} % prevent overfull lines
\providecommand{\tightlist}{%
  \setlength{\itemsep}{0pt}\setlength{\parskip}{0pt}}
\setcounter{secnumdepth}{5}
\ifLuaTeX
  \usepackage{selnolig}  % disable illegal ligatures
\fi
\IfFileExists{bookmark.sty}{\usepackage{bookmark}}{\usepackage{hyperref}}
\IfFileExists{xurl.sty}{\usepackage{xurl}}{} % add URL line breaks if available
\urlstyle{same}
\hypersetup{
  colorlinks=true,
  linkcolor={blue},
  filecolor={Maroon},
  citecolor={Blue},
  urlcolor={blue},
  pdfcreator={LaTeX via pandoc}}

\author{}
\date{}

\begin{document}

\hypertarget{ux6b21ux5143ux5b9aux7406ranknullity-theorem}{%
\section{次元定理（Rank--Nullity
Theorem）}\label{ux6b21ux5143ux5b9aux7406ranknullity-theorem}}

\texttt{A:R\^{}n→R\^{}m} に対して

\[
\boxed{\dim\ker A+\operatorname{rank}(A)=n}
\]

が成立します。

\hypertarget{ux672cux8ceaux7684ux306aux610fux5473}{%
\subsection{本質的な意味}\label{ux672cux8ceaux7684ux306aux610fux5473}}

入力空間の \texttt{n} 個の自由度は、

\begin{itemize}
\tightlist
\item
  変換後に残る方向
\item
  変換によって消える方向
\end{itemize}

に完全に分配されます。

例えば \texttt{rank(A)=r} なら、独立な出力情報は \texttt{r}
次元で、失われる入力自由度は \texttt{n-r} 次元です。

\hypertarget{svd-ux3067ux898bux308b}{%
\subsection{SVD で見る}\label{svd-ux3067ux898bux308b}}

\[
A=U\Sigma V^\top
\]

で非零特異値が \texttt{r} 個なら、\texttt{V} の最初の \texttt{r}
本は情報が出力へ伝わる方向、残り \texttt{n-r} 本は零空間方向です。

したがって次元定理は、SVD の幾何学ではほとんど目で見える形になります。

\end{document}
